\chapter*{Заключение}                       % Заголовок
\addcontentsline{toc}{chapter}{Заключение}  % Добавляем его в оглавление


В работе исследованы фазовые режимы спиновых систем на динамически изменяемой геометрии, задаваемой блужданиями без самопересечений на регулярных решётках, а усреднение наблюдаемых величин выполняется совместно по спиновым и геометрическим степеням свободы. 

Развит и реализован численный подход на основе метода Монте-Карло по  марковской цепи. 
Реализованный алгоритмический подход позволяет согласованно усреднять по спиновым и геометрическим степеням свободы и корректно выделять структурные и магнитные переходы в одном ансамбле.

Полученные и проанализированные результаты  компьютерного моделирования сводятся к следующему.

\begin{enumerate}
	 
	\item Для динамической HP-модели на квадратной решётке в большом каноническом ансамбле с использованием алгоритма Берретти--Сокала, обобщённого на случай динамических внутренних состояний мономеров, построена фазовая диаграмма и качественно исследовано геометрическое поведение цепи. Показано, что при малых значениях параметра взаимодействия модель ведёт себя подобно невзаимодействующему блужданию без самопересечений (показатель $\nu \to 3/4$). При усилении взаимодействия система переходит к более компактным конфигурациям, а область структурного перехода клубок-глобула смещена к большим значениям взаимодействия по сравнению с гомополимерным случаем, что обусловлено тем, что энергетический выигрыш возникает только для контактов типа $H$--$H$.  Результаты исследования динамической HP-модели послужили основанием для перехода к моделированию более сложных спин-полимерных систем в каноническом ансамбле.
	\item В двумерной модели Изинга на БСП геометрический переход ``клубок–глобула'' и магнитный переход имеют непрерывный характер; их локализация практически совпадает в пределах численных погрешностей. Формы распределений энергии остаются одномодальными, что согласуется с отсутствием латентной теплоты.
	\item В двумерной XY-модели на БСП структурный переход также непрерывен; магнитный переход является непрерывным и происходит при более сильном эффективном взаимодействии. Геометрический и магнитный критерии разделяются.
	\item В трёхмерной XY-модели на БСП структурный переход остаётся непрерывным, тогда как магнитный переход демонстрирует признаки первого рода (бимодальные распределения энергии и отрицательный провал кумулянта Биндера для намагниченности в узком диапазоне параметров).
	\item Геометрическое критическое поведение наследуется от ``родительской'' полимерной модели: добавление спинов не меняет класс универсальности \(\Theta\)-перехода для моделей с коротким действием.
	В то же время магнитное поведение существенно зависит от размерности системы и типа спиновых степеней свободы: в двумерном случае магнитные переходы остаются непрерывными (как для Изинга, так и для XY), тогда как в трёхмерной XY-модели возникает переход первого рода.
	  
\end{enumerate} 

\section*{Перспективы исследования}
 
 
 Реализованный вычислительный подход естественно можно адаптировать  к дальнодействию (например, эффективным потенциалам типа $J(r)\propto r^{-\alpha}$ или дипольно-подобным взаимодействиям), что позволит приблизить изучаемые системы к более реалистичным. 
 %Кроме того, развитие подхода в направлении моделей с большим числом спиновых состояний (модели Поттса и Блюма--Эмери--Гриффитса) и их сочетания с динамической HP-моделью может открыть новые возможности для исследования мультикритических явлений в спин-полимерных системах.
 
 
 
 
 